\begin{table}[H]
    \centering\caption{Kebutuhan Fungsional (FR) Sistem}\label{tbl:kebutuhan-fungsional}
% Menambah sedikit jarak antar baris agar tidak terlalu padat
\renewcommand{\arraystretch}{1.3}
\begin{tabular}{|l|p{0.85\textwidth}|}
\hline
\textbf{Kode} & \textbf{Deskripsi Kebutuhan Fungsional} \\
\hline

% Bagian 1
\multicolumn{2}{|l|}{\textbf{Interaksi Pengguna (pemelajar)}} \\
\hline
FR-01 & Sistem harus menyediakan mekanisme otentikasi bagi pemelajar untuk mengakses lingkungan pembelajaran personal mereka. \\
\hline
FR-02 & Sistem harus menyajikan konten pembelajaran, spesifiknya konten asesmen, yang telah disesuaikan tingkat kesulitannya dengan kemahiran pemelajar saat ini. \\
\hline
FR-03 & Sistem harus memfasilitasi \textbf{interaksi pembelajaran antar-pengguna} yang memungkinkan terjadinya pertukaran informasi atau kolaborasi akademik di dalam platform. \\
\hline
FR-04 & Sistem harus menyediakan informasi mengenai konteks sosial pembelajaran (misalnya, aktivitas komunitas atau capaian kolektif) untuk membangun kesadaran sosial (\textit{social awareness}). \\
\hline

% Bagian 2
\multicolumn{2}{|l|}{\textbf{Logika Adaptasi \& Pemrosesan Data (Sistem)}} \\
\hline
FR-05 & Sistem harus memiliki algoritma untuk mengestimasi tingkat kemahiran pengetahuan individu pemelajar secara dinamis berdasarkan riwayat aktivitasnya. \\
\hline
% FR-06 & Sistem harus mampu \textbf{mengkuantifikasi data interaksi sosial} (dari FR-03) menjadi parameter numerik yang dapat diproses lebih lanjut (misalnya, skor keterlibatan atau skor kinerja kolaboratif). \\
FR-06 & Sistem harus mampu mendeteksi kondisi stagnasi pembelajaran (\textit{learning plateau}) pemelajar secara otomatis berdasarkan suatu metrik atau proksi. \\
\hline
FR-07 & Sistem harus mampu merekomendasikan \textbf{intervensi sosial spesifik} (misalnya, merekomendasikan orang untuk berinteraksi atau topik untuk didiskusikan) agar bisa mendorong pemelajar keluar dari stagnasi pembelajaran. \\
\hline
FR-08 & Sistem harus menggunakan metrik kuantitatif untuk menentukan materi atau aktivitas berikutnya yang paling optimal bagi pemelajar. \\
\hline
FR-09 & Sistem harus memiliki mekanisme untuk menggabungkan alur pembelajaran yang adaptif secara individual (FR-05 dan FR-08) dengan aktivitas sosial (FR-07) untuk memberikan diversifikasi konten pembelajaran. \\
\hline
\end{tabular}
\end{table}